\begin{proofbox}
\textbf{方法一：代入焦点坐标，三类都只需一步（优先）。}
横向椭圆 $x^2/a^2+y^2/b^2=1$，右焦点 $F=(c,0)$，$c^2=a^2-b^2$。
通径所在直线为 $x=c$，代入得到
\[
 y^2=b^2\left(1-\frac{c^2}{a^2}\right)=\frac{b^4}{a^2},
 \qquad y=\pm\frac{b^2}{a}.
\]
因此两端距离为 $L=2b^2/a$。
横向双曲线 $x^2/a^2-y^2/b^2=1$ 中，$c^2=a^2+b^2$，同样代入 $x=c$：
\[
 y^2=b^2\left(\frac{c^2}{a^2}-1\right)=\frac{b^4}{a^2},
 \qquad L=\frac{2b^2}{a}.
\]
抛物线 $y^2=2px$ 的焦点横坐标是 $p/2$，故
\[
 y^2=2p\cdot(p/2)=p^2,\quad y=\pm p,\quad L=2p.
\]
参数均为正，开方得到的正长度分别是 $b^2/a$ 和 $p$。
对称、反射和交换坐标分别覆盖另一焦点、负向开口和纵向模型。
\end{proofbox}
\clearpage
\begin{proofbox}
\textbf{方法二：焦准比加对称性（几何角度）。}
设 $A,B$ 是横向椭圆或双曲线右通径端点，$F=(c,0)$，对应准线 $l:x=a^2/c$。
$AH\perp l$，且 $H$ 为垂足；通径竖直，所以
\[
 AH=\left|c-\frac{a^2}{c}\right|.
\]
\begin{center}\begin{tikzpicture}[font=\small,>=stealth]
\draw[black!30,rounded corners=2pt] (0,0) rectangle (14.4,8.2);
\node[anchor=west,font=\small\bfseries] at (.5,7.72) {几何证明：端点到焦点的距离，就是半通径};

\begin{scope}[shift={(5.8,4.1)},scale=.57]
\draw[blue,very thick] (0,0) ellipse (5 and 3);
\draw[black!35] (-5.5,0)--(7,0);
\draw[violet,thick] (6.25,-4)--(6.25,4);
\draw[orange!85!black,very thick] (4,-1.8)--(4,1.8);
\draw[black!65,thick] (4,1.8)--(6.25,1.8);
\fill (4,1.8) circle (3pt) node[above=12pt] {$A$};
\fill (4,-1.8) circle (3pt) node[below=12pt] {$B$};
\fill (4,0) circle (3pt) node[below left=6pt] {$F$};
\fill (6.25,1.8) circle (3pt) node[right=8pt] {$H$};
\node[left=7pt] at (4,.9) {$AF=9/5$};
\node at (0,4) {$AH=9/4$};
\end{scope}
\node at (9.3625,6.9) {$l:x=25/4$};
\node at (7.2,1.15) {$AF=(c/a)AH=(4/5)(9/4)=9/5$，$AB=18/5$};

\node[anchor=west] at (.7,.55) {$AH=|c-a^2/c|$，$AF=(c/a)AH=b^2/a$；上下对称后再乘 $2$};
\end{tikzpicture}\end{center}

由焦准比 $AF/AH=c/a$，
\[
 AF=\frac ca\left|c-\frac{a^2}{c}\right|
 =\frac{|c^2-a^2|}{a}=\frac{b^2}{a}.
\]
上下对称保证 $AF=BF$，故整弦 $AB=2AF=2b^2/a$。
这里的绝对值同时覆盖椭圆 $c<a$ 和双曲线 $c>a$。

\medskip
抛物线的通径所在直线与准线平行，两直线间距恰为焦准距 $p$。
端点到准线距离为 $p$，由等距定义，端点到焦点距离也为 $p$，再由对称性得 $L=2p$。
\begin{center}\begin{tikzpicture}[font=\small,>=stealth]
\draw[black!30,rounded corners=2pt] (0,0) rectangle (14.4,8.2);
\node[anchor=west,font=\small\bfseries] at (.5,7.72) {开口改变，只改变通径的位置和方向，不改变长度公式};
\begin{scope}[shift={(3.6,4.5)},scale=.38]
\draw[blue,very thick,domain=-5.5:5.5,samples=121,variable=\t] plot ({-\t*\t/8},\t);\draw[black!35] (-5.5,0)--(2,0);\draw[orange!85!black,very thick] (-2,-4)--(-2,4);\fill (-2,0) circle (3pt) node[below right=5pt] {$F$};\end{scope}
\node at (3.6,1.3) {$y^2=-8x$，$F=(-2,0)$};
\begin{scope}[shift={(10.8,4.5)},scale=.38]
\draw[blue,very thick,domain=-5.5:5.5,samples=121,variable=\t] plot (\t,{\t*\t/8});\draw[black!35] (0,-1)--(0,6);\draw[orange!85!black,very thick] (-4,2)--(4,2);\fill (0,2) circle (3pt) node[below right=5pt] {$F$};\end{scope}
\node at (10.8,1.3) {$x^2=8y$，$F=(0,2)$};

\node[anchor=west] at (.7,.55) {同尺度；两图 $p=4$，半通径 $4$，整条橙色弦都是 $2p=8$};
\end{tikzpicture}\end{center}

这种证明不依赖开口方向；平移、旋转后仍是同样的距离关系。
\end{proofbox}
\clearpage
